\documentclass{scrartcl}
\usepackage{evan}
%\usepackage[colorlinks=false, pdfborder={0 0 0}]{hyperref}

\begin{asydef}
  // Old cse5/olympiad.asy helpers this handout's diagrams still rely on,
  // reimplemented against stock Asymptote + evan.sty's own asydef.
  real markscalefactor = 0.03;

  // Right-angle mark at B for angle ABC.
  path rightanglemark(pair A, pair B, pair C, real s = 8) {
    pair P = s * markscalefactor * unit(A - B) + B;
    pair R = s * markscalefactor * unit(C - B) + B;
    pair Q = P + R - B;
    return P -- Q -- R;
  }

  // Ticks marking path g, matching the old pathticks().
  picture pathticks(path g, int n = 1, real r = .5, real spacing = 6,
      real s = 8, pen p = currentpen) {
    picture pict;
    real l = arclength(g);
    real space = spacing * markscalefactor;
    real halftick = s * markscalefactor / 2;
    if (n > 0) {
      pair direct = unit(dir(g, arctime(g, r * l)));
      real startpt = r * l - (n - 1) / 2 * space;
      for (int i = 0; i < n; ++i) {
        real t = startpt + i * space;
        pair B = point(g, arctime(g, t)) + (0, 1) * halftick * direct;
        pair C = B + 2 * (0, -1) * halftick * direct;
        draw(pict, B -- C, p);
      }
    }
    return pict;
  }

  // The point one unit from B along the internal bisector of angle ABC
  // (note: the vertex of the bisected angle is the *second* argument,
  // matching the old call sites).
  pair bisectorpoint(pair A, pair B, pair C) {
    pair P = unit(A - B) + B;
    pair R = unit(C - B) + B;
    return unit((P + R) / 2 - B) + B;
  }

  // The point one unit from the midpoint M of AB, along the perpendicular
  // bisector of AB.
  pair bisectorpoint(pair A, pair B) {
    pair M = midpoint(A -- B);
    return unit(rotate(90, M) * A - M) + M;
  }

  // The second point where circle c (assumed to pass through P) meets the
  // line through P and Q, other than P itself -- e.g. "let the
  // circumcircle meet side AB again at ...". Uses the analytic
  // circle-line intersection (geometry.asy), which is robust even when P
  // sits exactly on the circle, unlike intersecting c with the finite
  // segment P--Q via generic path intersection.
  pair secondIntersection(circle c, pair P, pair Q) {
    point[] pts = intersectionpoints(c, line(P, Q));
    return (abs(pts[0] - P) > abs(pts[1] - P)) ? (pair)pts[0] : (pair)pts[1];
  }

  pen pointpen = black;
  pen pathpen = black;
  pair D(pair p, pen q=pointpen) { return Drawing("", p, plain.S, q); }
  path D(path p, pen q=pathpen, arrowbar ar=None) { return Drawing(p, q, ar); }
  pair IntersectionPoint(path p, path q, int m=0) { return m == 0 ? IP(p, q) : OP(p, q); }
\end{asydef}

\ihead{Barycentric Coordinates}
\ohead{Berkeley Math Circle}
\ifoot{}
\cfoot{\thepage}


\begin{document}
\title{Barycentric Coordinates}
\author{Evan Chen \plusemail{evan@evanchen.cc}}
\date{October 23, 2012}

\maketitle

\begin{abstract}
  \noindent \footnotesize Much of this lecture is drawn from the significantly more extensive article \textit{Barycentric Coordinates in Olympiad Geometry}, which is \cite{bary} in the references.  It is a featured article on AoPS and can be found at \\ \url{http://www.artofproblemsolving.com/Resources/Papers/Bary_full.pdf}.
\end{abstract}

\section{Introduction}
Cartesian coordinates are great: they allow you to translate an arbitrarily complicated geometry problem into a straightforward (albeit tedious) computation.  However, the reference frame of Cartesian coordinates is a bit weird: you start with a pair of perpendicular lines and build everything from there.  This works fine for certain problems (USAMO 2012/5, anyone?), but is otherwise awkward.

\textit{Barycentric coordinates} follow the same theme, but are based off an arbitrary triangle, rather than a pair of perpendicular lines.  The three vertices of the triangle play the role of the origin, and the sides of the triangle play the role of the axes.

In this lecture, we will present these tools ``on the fly'' as they appear in a single problem.

\section{The Problem}
(Almost TSTST 2012) Triangle $ABC$ is inscribed in circle $\Omega$.  The interior angle bisector of angle $A$ intersects side $BC$ and $\Omega$ at $D$ and $L$ (other than $A$), respectively.  Let $M$ be the midpoint of side $BC$.  The circumcircle of triangle $ADM$ intersects sides $AB$ and $AC$ again at $Q$ and $P$ (other than $A$), respectively.  If $N$ is the midpoint of segment $PQ$, show that $MN \parallel AD$.  %Let $N$ be the midpoint of segment $PQ$, and let $H$ be the foot of the perpendicular from $L$ to line $ND$.  Prove that line $ML$ is tangent to the circumcircle of triangle $HMN$.

\begin{figure}[ht]
  \centering
  \begin{asy}
size(7cm); real lsf=0.9001; real lisf=2011.0; defaultpen(fontsize(8pt));

/* Initialize Objects */
pair A = (-5.0, 6.0);
pair B = (-6.0, -1.0);
pair C = (5.0, -1.0);
path c = A--B;
path a = B--C;
path b = C--A;
pair M = midpoint(a);
pair D = IntersectionPoint(Line(A,bisectorpoint(B,A,C),lisf),a);
circle w = circumcircle(A,D,M);
pair Q = secondIntersection(w,A,B);
pair P = secondIntersection(w,A,C);
pair L = IntersectionPoint(Line(A,D,lisf),Line(midpoint(relpoint(a, 0)--relpoint(a, 1)),bisectorpoint(relpoint(a, 0),relpoint(a, 1)),lisf));
pair N = midpoint(P--Q);
pair H = foot(L,N,D);

/* Draw objects */
dot(A);
dot(B);
dot(C);
draw(c, rgb(0.6,0.2,0.0));
draw(a, rgb(0.6,0.2,0.0));
draw(b, rgb(0.6,0.2,0.0));
dot(M);
dot(D);
draw(circumcircle(A,B,C));
draw(w);
dot(Q);
dot(P);
// dot(L);
draw(A--D);
// draw(D--L, dotted);
draw(P--Q);
dot(N);
draw(M--N);
// dot(H);
// draw(M--L, dotted);
// draw(H--N, dotted);
// draw(H--M, dotted);
// draw(H--L, dotted);

/* Label points */
label("$A$", A, lsf*dir(135));
label("$B$", B, lsf*dir(225));
label("$C$", C, lsf*dir(-15));
label("$M$", M, lsf*dir(75));
label("$D$", D, lsf*dir(45));
label("$Q$", Q, lsf*dir(45));
label("$P$", P, lsf*dir(90)*2.4);
// label("$L$", L, lsf*dir(-45)*2);
label("$N$", N, lsf*dir(45));
// label("$H$", H, lsf*dir(-115)*2);
\end{asy}
\caption{The Big Diagram}
\end{figure}

%For those of you with a lot of olympiad experience, it is an easy exercise to show that the problem is equivalent to proving that $MN \parallel AD$; thus, we can ignore the points $H$ and $L$ and focus on proving the parallel lines.
\eject



\section{Theory \& Technique}
Fix $\triangle ABC$, labelled counterclockwise, and let $a=BC$, $b=CA$, $c=AB$.  We'll let $[\mathcal P]$ denote the area of a polygon $\mathcal P$.

\subsection{Setting the Stage}
In barycentric coordinates, points are defined with three coordinates (instead of just two), but these coordinates must sum to $1$.  The point $P = (x,y,z)$ corresponds to the sum \[ \vec P = x \vec A + y \vec B + z \vec C. \]  Of course, we immediately have $A = (1,0,0)$, $B=(0,1,0)$ and $C=(0,0,1)$.  Because $x+y+z=1$, the choice of the zero vector doesn't matter.

In Cartesian coordinates, the equation of a line is of the form $Ax + By + C = 0$.  The analogous form in barycentric coordinates is as follows.
\begin{theorem}[Line]
  The equation of a line in barycentric coordinates is $ux + vy + wz = 0$, where $u,v,w$ are real numbers, unique up to scaling.
  \label{thm:line}
\end{theorem}
\begin{ques*}
  What does the line $z=0$ correspond to?
\end{ques*}

\begin{figure}[ht]
  \centering
  \begin{asy}
  size(8.01cm); real lsf=0.5; pathpen=linewidth(0.7); pointpen=black; real xmin=-4,xmax=4.01,ymin=-2.29,ymax=2.29;
  pen zzttqq=rgb(0.6,0.2,0), cqcqcq=rgb(0.75,0.75,0.75);
  pair A=(-1,2), B=(-2,-1), C=(2,-1), P=(-0.26,0.13), D=(0.19,-1), E=(0.43,0.57), F=(-1.43,0.71);
  D(A--B--C--cycle,zzttqq);
  fill(P--B--C--cycle, green);
  fill(P--C--A--cycle, cyan);
  fill(P--A--B--cycle, lightmagenta);
  D(A); MP("A",(-0.97,2.04),NE*lsf); D(B); MP("B",(-1.97,-0.96),NW*lsf); D(C); MP("C",(2.03,-0.96),NE*lsf); D(P); MP("P",(-0.23,0.17),dir(70)*lsf); D(D); MP("D",(0.22,-0.96),NE*lsf); D(E); MP("E",(0.46,0.61),NE*lsf); MP("F",(-1.4,0.76),NW*lsf); D(F);
  D(A--B,zzttqq); D(B--C,zzttqq); D(C--A,zzttqq); D(A--P); D(P--B); D(P--C); D(P--(0.43,0.57),linetype("4 4")); D(P--D,linetype("4 4")); D(P--F,linetype("4 4"));
  MP("x",incenter(P,B,C));
  MP("y",centroid(P,C,A));
  MP("z",incenter(P,A,B));
  \end{asy}
  \caption{Areas and barycentrics}
  \label{fig:area}
\end{figure}

There's a second (equivalent) definition we can use: if $P = (x,y,z)$, then $P$ is the point for which the areas\footnote{These areas are signed, so when $P$ lies on the opposite side as $A$ of $BC$, the area $[PBC]$ is considered negative.} $[PBC]$, $[PCA]$, and $[PAB]$ are in the ratio $x : y : z$.  For this reason, barycentric coordinates are sometimes called \textit{areal coordinates}.  See figure \ref{fig:area}.

%\begin{ques*}
%  In figure \ref{fig:area}, what are the coordinates of the point $D$?
%\end{ques*}

\subsection{The Centroid}

\begin{figure}[ht]
  \centering
  \begin{asy}
  size(6cm);
  pair A=(2.1,3.7), B=(0,0), C=(5.9,0);
  D(A--B--C--cycle);
  pair M_A = midpoint(B--C), M_B = midpoint(C--A), M_C = midpoint(A--B);
  D(A--M_A);
  D(B--M_B);
  D(C--M_C);
  pair G = centroid(A,B,C);
  D(A); D(B); D(C); D(M_A); D(M_B); D(M_C); D(G);
  add(pathticks(B--M_A, 1));
  add(pathticks(C--M_A, 1));
  add(pathticks(C--M_B, 2));
  add(pathticks(A--M_B, 2));
  add(pathticks(A--M_C, 3));
  add(pathticks(B--M_C, 3));

  MP("A", A, dir(45));
  MP("B", B, dir(135));
  MP("C", C, dir(45));
  MP("M_A", M_A, dir(45));
  MP("M_B", M_B, dir(45));
  MP("M_C", M_C, dir(135));
  MP("G", G, dir(65)*2);
\end{asy}
\caption{The centroid $G$ of $\triangle ABC$.  $AM_A$, $BM_B$ and $CM_C$ are medians.}
\end{figure}
The first thing we need to find in our original problem is the coordinates of the point $M$.  While we're here, we will also look at the other two midpoints, and the centroid, denoted $G$.

\begin{ques*}
  Determine the coordinates of the midpoints of $BC$, $CA$ and $AB$.
\end{ques*}
\begin{ques*}
  Determine the equation of the medians.
\end{ques*}
\begin{ques*}
  Determine the coordinates of the centroid $G$.
\end{ques*}
Notice that we managed to establish the existence of the point $G$ almost by accident; these same ideas can be modified to yield a proof of Ceva's Theorem.

\subsection{The Incenter}
The next point we need to find the coordinates of is the point $D$.  Along the way, we'll derive the angle bisector theorem.\footnote{Conversely, if you know the angle bisector theorem, you can find the coordinates of $D$ pretty quickly.  But shh!}

\begin{figure}[ht]
  \centering
  \begin{asy}
  size(6cm);
  pair A=(1.6,3.7), B=(0,0), C=(5.9,0);
  pair I=incenter(A,B,C);
  D(A--B--C--cycle);
  D(A--I);
  D(B--I);
  D(C--I);
  D(A); D(B); D(C); D(I);

  D(incircle(A,B,C), grey);
  D(I--foot(I,B,C), grey);
  D(I--foot(I,C,A), grey);
  D(I--foot(I,A,B), grey);
  D(rightanglemark(I,foot(I,B,C),C), grey);
  D(rightanglemark(I,foot(I,C,A),A), grey);
  D(rightanglemark(I,foot(I,A,B),B), grey);

  MP("A", A, dir(45));
  MP("B", B, dir(135));
  MP("C", C, dir(45));
  MP("I", I, dir(65)*2);
  \end{asy}
  \caption{The incenter $I$ of $\triangle ABC$.  $AI$, $BI$ and $CI$ are angle bisectors.}
\end{figure}

\begin{ques*}
  Show that $[IBC] : [ICA] : [IAB] = a : b : c$.
\end{ques*}
Using this, we can write $I = \left( \frac{a}{a+b+c} , \frac{b}{a+b+c} , \frac{c}{a+b+c} \right)$.  At this point, writing the denominators is getting cubersome, so we will abbreviate this as just $I = (a : b: c)$.  More generally, We will let $(x : y : z)$, where $x+y+z \neq 0$, denote the point $\left( \frac{x}{x+y+z}, \frac{y}{x+y+z}, \frac{z}{x+y+z} \right)$.

\begin{ques*}
  Can we plug $(a : b : c)$ instead of $\left( \frac{a}{a+b+c}, \frac{b}{a+b+c}, \frac{c}{a+b+c} \right)$ into the line formula (Theorem \ref{thm:line}) and still get valid results?
\end{ques*}
\begin{ques*}
  What are the coordinates of point $D$?
\end{ques*}
\begin{ques*}
  Deduce the angle bisector theorem from the coordinates of $D$.
\end{ques*}


\subsection{The Circle}
The next thing we need to do is find the equation of the circumcircle of $\triangle ADM$.  Let's call it $\omega$.

In Cartesian Coordinates, the general equation of a circle is $x^2 + y^2 + Ax + By + C = 0$.  The analog in barycentric coordinates is as follows:
\begin{theorem}[Circle]
  The equation of a circle in barycentric coordinates has the form \[ -a^2yz - b^2zx - c^2xy + (x+y+z)(ux+vy+wz) = 0 \]
\end{theorem}
Like its Cartesian analog, this isn't in general much fun to use.  Thankfully, in barycentric coordinates, several terms die if any of $\{x,y,z\}$ are zero.

\begin{ques*}
  Show that if a circle passes through $A=(1,0,0)$, then $u=0$.
\end{ques*}
\begin{ques*}
  Again, are we allowed to plug in $(a:b:c)$ instead of $\left( \frac{a}{a+b+c} , \frac{b}{a+b+c} , \frac{c}{a+b+c} \right)$ into the circle formula?
\end{ques*}

We can get the equation of $\omega$ just as we might in Cartesian coordinates.  We know from our earlier work that it passes through the points $A=(1:0:0)$, $M=(0:1:1)$ and $D=(0:b:c)$.  If we plug each of these in, we'll get three (linear) equations in three unknowns $u,v,w$

\begin{ques*}
  Show that $u=0$ and that \[ v = \frac{a^2c}{2(b+c)}, \quad w = \frac{a^2b}{2(b+c)}. \]
\end{ques*}

\subsection[The Points P and Q]{The Points $P$ and $Q$}
\begin{figure}[t]
  \centering
  \begin{asy}
size(6cm); real lsf=0.9001; real lisf=2011.0; defaultpen(fontsize(10pt));

/* Initialize Objects */
pair A = (-5.0, 6.0);
pair B = (-6.0, -1.0);
pair C = (5.0, -1.0);
path c = A--B;
path a = B--C;
path b = C--A;
pair M = midpoint(a);
pair D = IntersectionPoint(Line(A,bisectorpoint(B,A,C),lisf),a);
circle w = circumcircle(A,D,M);
pair Q = secondIntersection(w,A,B);
pair P = secondIntersection(w,A,C);
pair N = midpoint(P--Q);

/* Draw objects */
dot(A);
dot(B);
dot(C);
draw(c, rgb(0.6,0.2,0.0));
draw(a, rgb(0.6,0.2,0.0));
draw(b, rgb(0.6,0.2,0.0));
dot(M);
dot(D);
draw(circumcircle(A,B,C));
draw(w);
dot(Q);
dot(P);
draw(A--D);
draw(P--Q);
dot(N);
draw(M--N);

/* Label points */
label("$A$", A, lsf*dir(135));
label("$B$", B, lsf*dir(225));
label("$C$", C, lsf*dir(-15));
label("$M$", M, lsf*dir(75));
label("$D$", D, lsf*dir(45));
label("$Q$", Q, lsf*dir(45));
label("$P$", P, lsf*dir(90)*2.4);
label("$N$", N, lsf*dir(45));
\end{asy}
\caption{The Diagram Again}
\end{figure}
\begin{ques*}
  Find two constraints on the coordinates of $Q$, and use these to solve for $Q$.
\end{ques*}
\begin{ques*}
  Find the coordinates of $P$.  (Hint: use symmetry!)
\end{ques*}


\subsection[Finding N]{Finding $N$}
To take the midpoint of two points, simply normalize them, and take the average of their coordinates.  This is valid by the vector definition $x \vec A + y \vec B + z \vec C$.
\begin{ques*}
  Why is it important that the coordinates be normalized?
\end{ques*}
\begin{ques*}
  What is $N$?
\end{ques*}

\subsection{The Area Formula}
We're almost done now, because we've obtained the coordinates of all the points.  All that's left to do is show that $MN \parallel AD$.  But how do we do this?  We need one last tool, the area formula.\footnote{Can you show that this theorem implies the form for the equation of a line?}
\begin{theorem}[Area Formula]
  Let $P_i = \left( x_i, y_i, z_i \right)$ for $i=1,2,3$.  Then the area of $\triangle P_1P_2P_3$ is given by
  \[ [P_1P_2P_3] = [ABC] \cdot \left| \begin{array}{ccc} x_1 & y_1 & z_1 \\ x_2 & y_2 & z_2 \\ x_3 & y_3 & z_3 \end{array} \right| \]
  This area is positive when $P_1$, $P_2$ and $P_3$ are labelled counterclockwise, and negative otherwise.
  \label{thm:area}
\end{theorem}
\begin{ques*}
  Is it important that the $P_i$ are normalized?
\end{ques*}
\begin{ques*}
  We're trying to prove that $MN \parallel AD$.  What does this have to do with areas?
\end{ques*}
\begin{ques*}
  Show that the problem is solved upon verifying \[
  \left|
  \begin{array}{ccc}
    \frac{a^2}{4bc} & \half - \frac{a^2b}{4bc(b+c)} & \half-\frac{a^2c}{4bc(b+c)} \\
    0 & b & c \\
    1 & 0 & 0
  \end{array}
  \right|
  =
  \left|
  \begin{array}{ccc}
    0 & \half & \half \\
    0 & b & c \\
    1 & 0 & 0
  \end{array}
  \right|.
  \]
\end{ques*}

\section{Summary}
Here's a recap of the formulae that appeared:

\newcommand{\tableskip}{0.4cm}
\begin{table}[h]
  \centering
  \begin{tabular}{r|p{0.35\textwidth}|p{0.45\textwidth}}
    & Cartesian & Barycentric \\ \hline
    Definition & If $P=(x,y)$, then $\vec P = x \hat i + y \hat j$. & If $P=(x,y,z)$, where $x+y+z=1$, then $\vec P = x \vec A + y \vec B + z \vec C$. \\[\tableskip]
    Origin & Arbitrary point $(0,0)$. & Vertices of triangle $(1,0,0)$ etc. \\[\tableskip]
    Axes & Two perpendiculars $x,y=0$. & Sides of triangle $x,y,z=0$ \\[\tableskip]
    Line & $Ax + By + C = 0$ & $ux + vy + wz = 0$ \\[\tableskip]
    Circle & $x^2+y^2+Ax+By+C = 0$ & $-a^2yz-b^2zx-c^2xy \newline + (x+y+z)(ux+vy+wz)=0$ \\[\tableskip]
    Area & $\half \left|
    \begin{array}{ccc}
      x_1 & y_1 & 1 \\
      x_2 & y_2 & 1 \\
      x_3 & y_3 & 1
    \end{array}
    \right|$ &
    $[ABC] \cdot \left|
    \begin{array}{ccc}
      x_1 & y_1 & z_1 \\
      x_2 & y_2 & z_2 \\
      x_3 & y_3 & z_3 \\
    \end{array}
    \right|$
  \end{tabular}
  \caption{Cartesian vs. Barycentric}
  \label{tab:compare}
\end{table}

\eject
\section{Additional Problems}
\begin{enumerate}
  \item (Mongolia TST 2000/6) In a triangle $ABC$, the angle bisector at $A$,$B$,$C$ meet the opposite sides at $A_1$,$B_1$,$C_1$, respectively. Prove that if the quadrilateral $BA_1B_1C_1$ is cyclic, then \[ \frac{AC}{AB+BC} = \frac{AB}{AC+BC} + \frac{BC}{BA+AC}. \]
  \item (USA TST 2012) In acute triangle $ABC$, $\angle A < \angle B$ and $\angle A < \angle C$.  Let $P$ be a variable point on side $BC$.  Points $D$ and $E$ lie on sides $AB$ and $AC$, respectively, such that $BP = PD$ and $CP = PE$.  Prove that if $P$ moves along side $BC$, the circumcircle of triangle $ADE$ passes through a fixed point other than $A$.
   \item (USA TST 2003/2) Let $ABC$ be a triangle and let $P$ be a point in its interior. Lines $PA$, $PB$, $PC$ intersect sides $BC$, $CA$, $AB$ at $D$, $E$, $F$, respectively. Prove that \[ [PAF]+[PBD]+[PCE]=\frac{1}{2}[ABC]  \] if and only if $P$ lies on at least one of the medians of triangle $ABC$. (Here $[XYZ]$ denotes the area of triangle $XYZ$.)
  \item (Saudi Arabia TST 2012) Let $ABCD$ be a convex quadrilateral such that $AB = AC = BD$.  The lines $AC$ and $BD$ meet at point $O$, the circles $ABC$ and $ADO$ meet again at point $P$, and the lines $AP$ and $BC$ meet at point $Q$.  Show that $\angle COQ = \angle DOQ$.
\end{enumerate}

\begin{thebibliography}{99}
  \bibitem{bary} Max Schindler and C.  ``Barycentric Coordinates in Olympiad Geometry.'' \\ \url{http://www.artofproblemsolving.com/Resources/Papers/Bary_full.pdf}.
  \bibitem{abel} Zachary Abel.  ``Barycentric Coordinates.'' \\ \url{http://zacharyabel.com/papers/Barycentric\_A07.pdf}.
\end{thebibliography}

\end{document}
